% Post-CP22 candidate. All characteristic-two local steps are explicit.
\subsection{Three-torsion node permutations and gluing hearts}
\status{Proved and independently checked.}

\begin{lemma}[Finite eigenmaps for a three-torsion translation]
\label{lem:nonproduct-three-theta-finite-eigenmaps-candidate}
\textup{[Proved under the stated hypotheses; complete independent audits passed.]}
Let $(B,\Theta)$ be an indecomposable principally polarized abelian
surface over an algebraically closed field of characteristic two,
and let $0\ne g\in B[3]$. Each of the three eigenspaces of a lift
of $g$ acting on $H^0(B,3\Theta)$ has dimension three, is
basepoint-free, and defines a finite morphism $B\to\mathbf P^2$.
\end{lemma}
\begin{proof}
Put $E=B/\langle g\rangle$ and let $\pi:B\to E$ be the quotient.
The cyclic subgroup is isotropic for the theta pairing, so
$L=\mathcal O_B(3\Theta)$ descends to an ample line bundle $M$
on $E$. Its three possible descents differ by the three characters
of the cyclic cover; write them $M_i$. The direct-image character
decomposition identifies the three eigenspaces with $H^0(E,M_i)$.
Riemann--Roch gives $M_i^2=6$ and $\chi(M_i)=3$; ample vanishing
gives $h^0(M_i)=3$. These standard cohomology formulas and their
characteristic-independent proofs are
\citet[Theorems~9.11 and~9.14, Proposition~9.18]{EdixhovenVanDerGeerMoonenAV}.

We give the basepoint argument, including fixed divisors. First no
elliptic curve $F\subset E$ has $M_i\cdot F=1$. Translate $F$ to
an elliptic subgroup and let $F'$ be the identity component of its
inverse image in $B$. The restriction $F'\to F$ has degree
$e\in\{1,3\}$. The projection formula gives
$3\Theta\cdot F'=e$, hence $e=3$ and $\Theta\cdot F'=1$.
An elliptic curve of principal degree one splits a principally
polarized surface as a product, contrary to the hypothesis.

Suppose the complete system $|M_i|$ has a nonzero fixed divisor
$D$, and write $M_i\sim D+N$, with $N$ the moving part.
Every effective divisor on an abelian variety is nef: translate
it away from any prescribed integral curve to compute a
nonnegative intersection. Thus both $D,N$ are nef.
If $N^2>0$, then $N$ is ample on the surface. Indeed the Hodge
index theorem shows that a curve with $N$-intersection zero
would have negative square, impossible for an effective divisor
on an abelian surface; Nakai--Moishezon then applies.
Ample vanishing and $h^0(N)=h^0(M_i)=3$ give $N^2=6$.
But $(M_i-N)\cdot(M_i+N)>0$, since $D\ne0$, $M_i$ is ample,
and $N$ is nef, contradicting $M_i^2=N^2=6$.

If $N^2=0$, each irreducible component has square zero and any
two have intersection zero. By adjunction its normalization has
genus at most one; an abelian variety contains no rational curve,
so the components are translates of elliptic subgroups. Distinct
elliptic directions have positive intersection, hence they are
all fibres of one elliptic quotient. Thus $N$ is the pullback of
a degree-$m$ divisor on the elliptic quotient and $h^0(N)=m$.
It follows that $m=3$. Writing its fibre class as $F$, we obtain
\[
              6=M_i^2=D^2+6D\cdot F.
\]
Ampleness makes $D\cdot F\ge1$, while $D^2\ge0$, so
$D\cdot F=1$ and $M_i\cdot F=1$, again impossible.
Consequently $|M_i|$ has no fixed divisor.

Its remaining base locus is finite. The theta group translates
it by $K(M_i)$, an \etale group of order $\chi(M_i)^2=9$.
Translations act freely on the surface, so a base point would
give at least nine distinct base points. Two general independent
members, having no common curve component, intersect with total
length $M_i^2=6$. This contradiction proves global generation.
The associated morphism $E\to\mathbf P^2$ is finite, since its
pullback of $\mathcal O(1)$ is the ample bundle $M_i$; its
composition with $\pi$ is finite as well.
\end{proof}

\begin{lemma}[The theta difference cover and its independent shifted branches]
\label{lem:theta-difference-cover-char2-candidate}
\textup{[Proved under the stated hypotheses; complete independent audits passed.]}
Let $C=\Theta\subset B$ be a smooth symmetric principal theta
curve in characteristic two. Away from zero the difference map
\[
            \beta:C\times C\longrightarrow B,\qquad(P,Q)\mapsto P-Q
\]
is finite flat of degree two and generically \etale. Its only
divisorial branch support is the reduced curve
$\Delta=[2]C$, of numerical class $2\Theta$.
The inertia permutation at its generic point is the nontrivial
transposition. No nonzero odd torsion translation preserves
$\Delta$.

For three theta translates with translating points
$a_1+a_2+a_3=0$, the two-point intersection covers for any one
pair and any prescribed shift by an odd torsion point have an
inertia transposition which fixes all covers belonging to other
pair/shift branch divisors.
\end{lemma}
\begin{proof}
For a nonzero difference, the fibre is the intersection of two
distinct theta translates and has length two. Properness gives
finiteness away from zero; the source is Cohen--Macaulay and the
target is smooth of the same dimension, giving flatness.
The differential fails to be invertible exactly when the two
translated tangent lines agree. These are the canonical tangent
directions of the genus-two curve. Apart from the diagonal, whose
image is zero, equality means $Q=\iota P$. Thus the only
divisorial branch support is $[2]C$.
The map is generically \etale since the canonical map is
nonconstant. Its quadratic generic extension is Galois, with
involution $(P,Q)\mapsto(\iota Q,\iota P)$.

By Lemma~\ref{lem:all-theta-double-image-candidate}, the ramification
curve $Q=\iota P$ maps to its reduced image with inseparable
degree two and separable degree one, and that image has class
$2\Theta$. In particular the quadratic cover is not \etale at
the generic point of this divisor. Over its strict henselian
local base its nontrivial separable generic extension therefore
has inertia of order two. This includes the possible purely
inseparable residue extension; no tame-ramification assertion
is made in characteristic two.

If translation by $h$ preserves $\Delta$, it preserves its line
bundle, hence $h\in K(\mathcal O(\Delta))=B[2]$; numerical
equivalence suffices here because the polarization homomorphism
depends only on the N\'eron--Severi class. An odd torsion $h$ is
therefore zero. Distinct odd shifts of $\Delta$ are distinct
irreducible divisors.

Finally the map from $(a_1,a_2)$, with $a_3=-a_1-a_2$, to any
two independent pair differences is an isogeny with determinant
three, hence \etale in characteristic two. Branch divisors
for a fixed pair are distinct by the preceding translation
argument. Those for different pairs are distinct because the
two corresponding differences vary independently. At a generic
point of any chosen branch divisor all the other covers are
\etale. Its inertia consequently acts as the one specified
transposition and as the identity on the other covers.
\end{proof}

\begin{proposition}[The full product of eighteen-node permutation groups]
\label{prop:nonordinary-three-torsion-node-product-candidate}
\textup{[Proved under the stated hypotheses; complete independent audits passed.]}
Let $(B,\Theta)$ be nonordinary and indecomposable, and let
$K\subset B[3]$ be a subgroup of order nine defining a polarized
quotient $\phi:B\to A$. On the generic norm family of curves in
$|3\Theta|$, the additional nodes have permutation image
\[
                      \prod_{h\in(K\setminus0)/\pm} S_{18}.
\]
\end{proposition}
\begin{proof}
Fix $0\ne g\in K$. Its eighteen oriented node points are
$P$ satisfying $s(P)=s(P+g)=0$; this is in bijection with the
eighteen nodes of unordered type $\{g,-g\}$ when the norm is
nodal. The one-point incidence is a projective-six bundle over
$B$, hence irreducible. We prove two-transitivity by examining
two distinct such points $P,Q$.

Choose the eigenbasis in the first lemma. An evaluation vector
has three nonzero components $v_0(P),v_1(P),v_2(P)$, and the
orbit plane is their direct sum. The chord
$\langle v(P),v(P+g)\rangle$ contains no eigen-coordinate vertex:
canceling two different eigenvalues in a combination of its
two spanning vectors would cancel the whole combination.
If the two chords for $P,Q$ meet, their orbit planes must
therefore meet in at least two coordinate directions. Thus at
least two of the eigen-projections of $P,Q$ coincide. Each
eigenmap is finite, so the locus of such pairs has dimension
at most two. Outside it the four value conditions are
independent. The rank-four open in $B\times B$ has dimension
four and is irreducible; its coefficient incidence is a
projective-four bundle of total dimension eight.

At four distinct points the rank is at least three, since any
three values of $3\Theta$ are independent. The exceptional
pair locus just bounded therefore contributes dimension at
most $2+(8-3)=7$. If $Q=P+g$ or $P-g$, only three distinct
points occur, with three independent values and a two-dimensional
point locus, giving the same bound. Hence these loci cannot
contribute a component of the finite \etale ordered-two-point
cover over the generic coefficient open. That cover is
irreducible, proving two-transitivity.

The needed transpositions can be obtained on a three-theta
boundary without imposing an unproved structured higher-jet
condition. By
Corollary~\ref{cor:nonordinary-theta-three-torsion-nodal-norm-candidate},
each individual theta norm has only its eight distinct nodes.
Choose three general translates with translating points of sum
zero. Their norms have distinct self-nodes, mutually transverse
intersections, and no triple points. To check the last assertion
directly, the map
\[
 (P,a_1,a_2)\longmapsto(P-\phi a_1,\ P-\phi a_2,\ P+\phi a_1+\phi a_2)
\]
is a finite \etale isogeny of three surfaces (its scalar
determinant is three). The inverse image of the product of
the three norm curves has dimension three, while the translating
parameter space has dimension four. Pair tangencies and incidence
with the finitely many self-nodes are likewise proper conditions,
by generic separability of the difference map.

In a fixed node type, six points are intrinsic self-nodes of the
three translated theta curves and are constant. The remaining
twelve form six two-point covers: the two orientations of each
of the three component pairs. Their branch supports are exactly
the shifted difference divisors of the preceding lemma. A
generic point of one such divisor supplies a transposition in
this type, fixing every other type and every other node.
These inertia elements belong to the generic norm-node image:
the finite \etale additional-node cover extends over the normal
coefficient open containing the generic three-theta boundary;
its smooth-source dense open surjects on the relevant geometric
fundamental group. Two-transitivity and one transposition give
$S_{18}$ in each projection. Normal conjugation of an isolated
factor transposition by this surjective projection yields the
whole $S_{18}$ supported in that factor. The full product follows.
\end{proof}

\begin{proposition}[All type hearts occur as gluing translations]
\label{prop:nonordinary-three-torsion-gluing-hearts-candidate}
\textup{[Proved under the stated hypotheses; complete independent audits passed.]}
In the setting of the preceding proposition, write the relative
norm sequence on its smooth maximal-rank open as
\[
                    0\longrightarrow T\longrightarrow V
                       \longrightarrow U\longrightarrow0,
       \qquad T=\bigoplus_h T_h,\quad \dim T_h=18,\quad\dim U=8.
\]
Its monodromy contains every translation in
$\Hom(U,T_h^0)$, where $T_h^0$ is the coordinate-sum-zero
hyperplane, for every type $h$.
\end{proposition}
\begin{proof}
Let $N=\prod_hS_{18}$. The joint image on $T$ and $U$ is
$N\times\GL_8(\mathbf F_2)$: each projection is surjective,
and the simple group $\GL_8(\mathbf F_2)$ has the prime divisor
31, which does not divide $|N|$, so Goursat's lemma excludes
a common nontrivial quotient. The kernel of this joint image
inside $\Hom(U,T)$ is $U^*\otimes T'$ for an $N$-stable
subspace $T'\subset T$. This is the elementary full-matrix-algebra
consequence for the natural $\GL_8$ module used in the actual
block proof. Modulo $T'$ the translation kernel is zero.
The vanishing $H^1(\GL_8,U^*)=0$ conjugates the restriction
of this extension to $\GL_8$ to a block-diagonal action.
The remaining $N$ factor commutes with $\GL_8$; its possible
upper-right column is therefore a $\GL_8$-invariant element
of $U^*\otimes(T/T')$, and is zero. We obtain an equivariant
splitting
\[
                         V/T'=(T/T')\oplus U.             \tag{NTG1}
\]

We justify using this splitting at the smaller-rank three-theta
boundary. Let $\mathcal B$ be the normal open of the source
projective coefficient space where all additional norm
identifications are distinct transverse pairs of smooth source
points. It contains the generic three-theta members; source
internal nodes and their selected-pair tangencies are allowed.
The partial-normalization structure-sheaf sequence gives on
$\mathcal B$ an exact sequence of coherent Frobenius bundles
\[
                         0\to T\otimes\mathcal O
                              \to\mathcal V\to\mathcal U\to0.
\]
The ambient $H^1(\mathcal O)$ maps are identified by the odd
isogeny. Both coherent relative bundles have constant ranks,
although their unit parts may drop. The graph term is a
Frobenius-unit bundle associated with the finite \etale
permutation system $T$.

The subspace $T'$ extends to a finite \etale subgroup of $T$
on $\mathcal B$: the dense maximal-rank open surjects on
$\pi_1(\mathcal B)$, and its action on $T$ already factors
through that group. On Lang schemes the projection
\[
              \mathcal L_{\mathcal V}/T'
                         \longrightarrow\mathcal L_{\mathcal U}
\]
is a finite \etale torsor under $T/T'$. This can be checked
after trivializing the graph bundle, when the equations for
a lift are $t_i^2-t_i=b_i(u)$.
The generic splitting (NTG1) is a section of this torsor on
the dense maximal-rank open. The full Lang scheme
$\mathcal L_{\mathcal U}$ is normal, being \etale over
the normal base. Every component dominates the base. A
section of a finite \etale morphism over a dense open of a
normal integral component extends uniquely: its closure is
finite and birational and hence is that normal component.
Thus (NTG1) extends to every finite Lang part at the boundary,
compatibly with addition and all monodromy. No full-rank
specialization isomorphism has been asserted.

Now work in the three-theta translating family and take inertia
at the branch divisor of a selected, unshifted component-pair
intersection. The difference-cover lemma makes this inertia
swap exactly its two selected nodes and fix all additional
norm nodes. Its action on cohomology of the nodal union has
rank one, with image
\[
                     q_{ij}=e_{ij,1}+e_{ij,2}
\]
in graph cohomology. Here is the local gluing justification in
characteristic two. The component curves are constant after
translation, so their finite \etale torsors are constant. Near
the coalescing pair, each such torsor is \etale and is trivial
over the strict henselian local branch neighbourhood. Its two
gluing labels can therefore be chosen in the same local
trivialization. Inertia exchanges these two labels and fixes
all other labels. Hence its image is exactly the displayed
one-dimensional graph line; a possible change of splitting
does not alter that image. This argument applies also when
the residue extension of the quadratic collision is purely
inseparable.

On the source doubled-triangle graph all three classes
$q_{12},q_{13},q_{23}$ coincide with one nonzero class $q$,
because their pairwise sums are vertex coboundaries. This
graph class remains nonzero in relative cohomology, since
ambient $H^1(B,\mathbf F_2)$ injects into each theta
normalization and meets the graph part trivially.
In (NTG1) the additional graph factor is fixed by each of
these inertia elements. The common source image $q$ must
therefore have the same lift in $V/T'$ for all three pairs.
Consequently
\[
                              q_{12}+q_{13}\in T'.       \tag{NTG2}
\]
In the target graph, this sum is the sum of all unselected
edges incident with vertex one, since the full vertex star
is a coboundary. In each type its eighteen coordinates have
eight ones on incident cross edges and ten zeros on the
other cross edges and the six loops. It is nonconstant in
every type. A transposition in any one $S_{18}$ factor
subtracts two unequal coordinates in that factor only;
normal conjugation then gives its entire sum-zero heart.
Thus $T_h^0\subset T'$ for every $h$, which is precisely
the claimed gluing-translation subgroup.
\end{proof}

\begin{theorem}[The degree-three image from a nonproduct three-neighbour]
\label{thm:nonordinary-degree-three-from-neighbour-candidate}
\textup{[Proved under the stated hypotheses; complete independent audits passed.]}
Let $\phi:(B,\Theta_B)\to(A,\Theta_A)$ be an \etale polarized
isogeny of multiplier three and degree nine in characteristic two,
with $B$ nonordinary and indecomposable. Then the relative
degree-three geometric monodromy on the smooth maximal-rank
open of $|9\Theta_A|$ is $\GL_{80}(\mathbf F_2)$.
The arithmetic image over any perfect field of definition is
the same full group.
\end{theorem}
\begin{proof}
The individual theta norms are nodal by
Corollary~\ref{cor:nonordinary-theta-three-torsion-nodal-norm-candidate}.
The three general translates and all their kernel-translated
partial normalizations satisfy the first geometric hypothesis
of Theorem~\ref{thm:three-theta-conditional-full-image-candidate},
by the incidence argument in the node-product proposition.
All the source unions whose translating points have sum zero
belong to the same $|3\Theta_B|$ and to the same connected
coefficient open. The preceding two propositions establish
the full type permutation groups and every gluing heart.
The actual natural $\GL_8$ block is already supplied by
Theorem~\ref{thm:nonordinary-ppav-odd-natural-block-candidate}.
Thus both geometric hypotheses of the conditional theorem
hold. Applying it with $d=3$ proves the geometric image
$\GL_{9\cdot3^2-1}=\GL_{80}$. All constructions may be
made after geometric base change; the arithmetic image
contains this full geometric group in the same general
linear group, so equality follows over the stated base.
\end{proof}

\begin{remark}[Remaining scope]
\label{rem:nonordinary-degree-three-neighbour-scope-candidate}
\textup{[Proved under the stated hypotheses; complete independent audits passed.]}
The theorem as stated assumes a nonproduct polarized three-neighbour.
Existence of such a neighbour for every fixed nonproduct nonordinary
surface is a separate assertion. The eigenmap proof uses that the
kernel is killed by three; neither it nor the theta nodality proof
has been extended here to arbitrary odd multiplier.

The finite graph supplement
\path{agents/monodromy/code/degree_three_theta_graph_checks.py}, with
full output
\path{agents/monodromy/outputs/degree_three_theta_graph_checks.json},
checks all 81 sum-zero translating triples in $K=\mathbf F_3^2$.
It reproduces graph dimension 76, partial-normalization kernel
dimension 72, the four-heart dimension 68, the total heart span
dimension 75, the rank-four constraints on type-constant vectors,
and the eight-one/ten-zero vertex star in each type. The author
run at 19:41:14.217748 UTC on 19 November 2026 passed in
0.012680367006396409 seconds. The parent and reduction/descent agent
independently read every proof and the complete program, and reran
the latter without changing the author output in respectively
0.014382761997694615 and 0.013090728003589902 seconds; both passed.
This finite supplement does not
replace any geometric incidence, inertia, or gluing argument.
\end{remark}
